\documentclass[10pt,a4paper]{amsart}
\usepackage[margin=19mm]{geometry}
\usepackage{amsmath,amssymb,mathtools}
\usepackage[scr=rsfs,cal=euler]{mathalfa}
\usepackage{xfrac}
\usepackage[unicode,hidelinks,bookmarks=false]{hyperref}
\newcommand{\ie}{i.e.\ }
\newcommand{\cf}{cf.\ }
\newcommand{\resp}{respectively\ }
\newcommand{\Subsection}{\subsection}
\providecommand{\nobreakdash}{\nobreak\hbox{-}}
\setlength{\emergencystretch}{2em}
\multlinegap0pt
\usepackage{amssymb}
\usepackage{xcolor}
\newtheorem{lemma}{Lemma}
\newcommand\Hom{{\rm Hom}}
\newcommand\Isom{{\mathbf{Isom}}}
\newcommand\Lspace{\mathsf E}
\newcommand\PO{{\mathsf{PSO}}}
\newcommand\SO{{\mathsf{SO}}}
\newcommand\bR{{\mathbb{R}}}
\newcommand\bZ{{\mathbb Z}}
\newcommand\ra{\rightarrow}
  \newcommand{\rome}{\mathrm e}
   \newcommand{\romel}{{\mathrm{e}, \mathrm{l}}}
\title{The stability of Margulis space-times with parabolic holonomy elements}
\author{Suhyoung Choi}
\date{Source: arXiv:2606.11694v1 (CC BY 4.0)}
\begin{document}
\maketitle

\noindent\emph{Source and licence.} The stability of Margulis space-times with parabolic holonomy elements. Version arXiv:2606.11694v1 (CC BY 4.0), distributed by arXiv under the Creative Commons Attribution 4.0 International licence, http://creativecommons.org/licenses/by/4.0/, as recorded in the arXiv licence metadata for this exact version and shown on the abstract page https://arxiv.org/abs/2606.11694v1. Full licence notice: \texttt{licenses/arxiv-2606.11694-CC-BY-4.0.txt}.\par\smallskip

\noindent\textit{Editorial scope.} Counted unit: the article's lemma lem:deformnop (source0.tex lines 1282--1295). The article's complete original proof of it is delivered with it (source0.tex lines 1296--1322, one \texttt{proof} environment of 27 lines). No mathematics is written, repaired or re-derived: the statement and the proof are the article's own lines, in the article's own notation, and no retained line is substituted or re-flowed.  No further source-local context is needed: every cross-reference of every retained line resolves inside the two delivered spans.  Nothing was omitted from any retained span, so the package carries no omission record.  No retained line contains a citation, so the wrapper carries no bibliography and no bibliography entry.  No retained line contains a figure environment, an image-inclusion command or drawing code, so no image is delivered and the package declares no asset dependency.  Editorial formatting only, in addition: an AMS \texttt{amsart} preamble that restates the article's own declarations and macros used by the retained lines, namely the environments \texttt{lemma} on separate counters, together with the standard packages those lines need.  The retained environments print in this document as Lemma 1 in order of appearance, whereas the article prints the same items with the numbering of its own full document.  The complete native source of the article as distributed in the arXiv e-print for this exact version is bundled unchanged as \texttt{sources/202568/source0.tex}.
\begin{lemma} \label{lem:deformnop} 
Let $\rho:F_n \ra \Isom(\Lspace)$ be an affine deformation of a Fuchsian representation whose image contains parabolic elements.
%where the image of 
%$\mathcal{L}\circ \rho$ is a free group with a parabolic element. 
Then $\mathcal{F}^{\rome}$ and $\mathcal{F}^{\romel}$ are semialgebraic sets. 
There is a cone neighborhood $N$ of $\rho$ in $\mathcal{F}^{\rome}_n$ where 
the subset $N'$ of $N$ consisting of representations without parabolic image elements is 
an open dense cone, which is a complement of a semialgebraic subset of codimension 
$\geq 1$ in $N$ passing through $\rho$. 
%Then the subset of elements $x$ of $D'_n$
%$\Gamma_x$ has no parabolic is an open dense subset of $D'_n$, which is the complement of finite union of codimension $1$ semi-algebraic set. 
%there is a path of affine deformations of Fuchsian representations $\rho_t$, $t \in [0, 1]$ where,
%$\rho_0 = \rho$ and $\rho_t$ has no parabolic for $t> 0$. 
\end{lemma}

\begin{proof} 
%Let $\mathcal{F}^{\rome}_n \subset \Hom(F_n, \SO(2,1))$ is the inverse image of $\mathcal{F}_n$. 
The subset $\hat D_n$ of $\mathcal{F}^{\rome}_n$ of representations 
with some parabolic image elements is a union of finitely many semi-algebraic subsets of codimension $\geq 1$ passing through $\rho$
since only peripheral elements can have a parabolic holonomy. 
The subset in question is the inverse image of $\hat D_n$. 

%There is a map 
%$\mathcal{L}': \Hom(F_n, \Isom(\Lspace)) \ra \Hom(F_n, \PO(2, 1))$ by %%sending 
%$h(g)$ to $\Pi\circ \mathcal{L}(h(g))$ for $g \in F_n$. 
Define $\mathcal{L'} := \Pi\circ \mathcal{L}\circ h$.
This can be considered a fiber bundle
$(\Pi\circ\mathcal{L})^n:\Isom(\Lspace)^n \ra \PO(2, 1)^n$
with fibers that can be identified with $(\bZ_2\ltimes \bR^{2, 1})^n$. 
The proposition follows from the standard 
cone-neighborhood theory of semi-algebraic sets. 
%Let $\mathcal{F}$ denote the subspace of Fuchsian representations
%in $ \Hom(F_n, \PO(2, 1))$. This fibers over a Fricke space of hyperbolic surfaces 
%allowing cusps, which is a cell with nonempty boundary.
%
%Since it is a fiber bundle, 
%$\mathcal{F}$ is a manifold with boundary. 
%We have $h \in ({\mathcal{L}^n})^{-1}(\mathcal{F})$, and 
%$({\mathcal{L}^n})^{-1}(\mathcal{F}) \ra \mathcal{F}$ is again a fiber bundle with fibers diffeomorphic to  $(\bR^{2, 1})^n$. 
%%Hence, there is always a path of Fuchsian representation $\rho_t$ where 
%$\rho_t$ has no parabolics. 
\end{proof} 

\end{document}
